\documentclass[11pt,a4paper,twocolumn]{article}

% --- PAQUETES ESTÁNDARES Y SEGUROS ---
\usepackage[utf8]{inputenc}
\usepackage[spanish, es-tabla]{babel} % Configura idioma y renombra "Cuadro" a "Tabla"
\usepackage[margin=2.3cm, headheight=15pt, headsep=0.7cm]{geometry} % Márgenes del documento
\usepackage{amsmath, amssymb, amsfonts} % Herramientas matemáticas para econometría
\usepackage{booktabs} % Diseña tablas limpias estilo Journal (estilo AEA)
\usepackage{tabularx} % Tablas que se ajustan automáticamente al 100% del ancho de página (\textwidth)
\usepackage{caption} % Personalización de títulos de tablas y figuras
\usepackage{fancyhdr} % Control de encabezados profesionales
\usepackage{setspace} % Control de interlineados

\usepackage{mathptmx} % Fuente Times New Roman (Texto y Matemáticas) 100% nativa y segura para MiKTeX

% --- CITACIÓN APA 7 Y ENLACES INTERACTIVOS (VINCULACIÓN) ---
\usepackage{xcolor} % Paquete de gestión de colores
\definecolor{darkblue}{RGB}{0, 45, 135} % Azul oscuro elegante y seguro para los enlaces
\usepackage[round, authoryear]{natbib} % Permite formato (Autor, Año) estilo APA 7
\usepackage{url} % Gestión de URLs
\usepackage[colorlinks=true, linkcolor=darkblue, citecolor=darkblue, urlcolor=darkblue, breaklinks=true]{hyperref} % Enlaces cliqueables en tablas y referencias
\makeatletter
\g@addto@macro{\UrlBreaks}{\UrlOrds}
\makeatother

% --- CONFIGURACIÓN DE COLUMNAS ---
\setlength{\columnsep}{0.75cm} % Espaciado elegante entre las dos columnas

% --- CONFIGURACIÓN DE ENCABEZADOS Y PIE DE PÁGINA ---
\pagestyle{fancy}
\fancyhf{} 
\fancyhead[L]{\footnotesize\textit{Canchi, A. et al.}} 
\fancyhead[R]{\footnotesize\textit{Semillas certificadas y rendimiento agrícola en el Perú 2024}} 
\fancyfoot[C]{\small\thepage} % Número de página ABAJO Y CENTRADO a partir de la página 2
\renewcommand{\headrulewidth}{0.4pt} % Línea divisoria superior sutil como en un journal real
\renewcommand{\footrulewidth}{0pt} % Sin línea inferior

% Estilo exclusivo para la página 1 (portada totalmente limpia, sin número de página ni línea)
\fancypagestyle{plain}{
  \fancyhf{} % Limpia encabezados y número de página en la página 1
  \renewcommand{\headrulewidth}{0pt} % Sin línea superior en la portada
}

% --- CAPTION ESTILO ELSEVIER / APA (Título en negrita, salto de línea, alineado a la izquierda sin cursiva) ---
\captionsetup[table]{labelfont=bf, textfont=normalfont, labelsep=newline, justification=raggedright, singlelinecheck=false, position=top}

\begin{document}

% --- PORTADA Y ABSTRACT EN UNA COLUMNA ---
\twocolumn[
  \begin{center}
    \vspace*{3.5cm} 
    
    {\LARGE \bfseries Efecto del uso de semillas certificadas sobre el rendimiento agrícola del cultivo de arroz cáscara en las unidades agropecuarias del Perú: Evidencia del año 2024 \par}
    \vspace{0.8cm}
    
    {\normalsize Ccanchi, A.; Condori, A.; Flores, M.; López, F.; Mamani, E.; Medina, C.; Morante, C.; Quispe, R.; Suri, D.; Villavicencio, S.; Yauri, G. \par}
    \vspace{0.6cm}
    
    \begin{minipage}{0.88\linewidth}
      \small
      \vspace{0.3cm}
      \begin{center}
        \textbf{\large Abstract}
      \end{center}
      \vspace{-0.15cm}
      This study evaluates the empirical impact of certified seed adoption on agricultural yield among rice farmers in Peru during the 2024 cropping season, drawing methodological foundation from Takeshima et al. (2025). Using a refined national microdata sample of $1{,}991$ agricultural units from the National Agricultural Survey (ENA 2024 - INEI), we estimate an Ordinary Least Squares (OLS) multiple regression model with Huber-White heteroscedasticity-robust standard errors and departmental fixed effects. Controlling for water access, climatic shocks, and sociodemographic profiles ($R^2 = 0.5480$, $F = 76.32$), the empirical findings validate our central hypothesis: certified seed adoption directly and significantly increases rice physical yield by $+9.82\%$ ($p < 0.001$) compared to common or recycled grain. This productivity gain is driven by the superior genetic purity, physiological vigor, and phytosanitary quality of certified germoplasma. Furthermore, technified water access yields a major boost of $+48.58\%$, whereas adverse climatic shocks and biological pests impose penalties between $-8.18\%$ and $-36.29\%$. We conclude that varietal certification and irrigation infrastructure constitute complementary, strategic policy determinants for boosting national agricultural productivity.

      \vspace{0.18cm}
      \noindent \textbf{Keywords:} certified seeds, agricultural yield, rice, ENA 2024, technified irrigation, climate shocks, robust OLS.
    \end{minipage}
    \vspace*{0.5cm}
  \end{center}
]

\thispagestyle{empty} 
\clearpage 

\section{Introducción}
La seguridad alimentaria y la productividad agrícola son pilares esenciales para el desarrollo socioeconómico actual. Ante el crecimiento demográfico, la adopción de semillas certificadas se ha consolidado como una estrategia eficiente para incrementar el rendimiento de los cultivos básicos sin expandir desmedidamente la frontera agrícola.

El arroz es fundamental para la seguridad alimentaria y la economía de las unidades agropecuarias peruanas, con producción concentrada en la costa y selva. No obstante, el sector enfrenta una brecha tecnológica debido a la adopción heterogénea de insumos y la persistencia del uso de semillas informales o recicladas, lo que limita la competitividad frente a quienes utilizan semillas certificadas.

Frente a este escenario, surge la necesidad de evaluar el impacto de esta innovación tecnológica en el campo peruano, es por ello que el presente estudio se plantea la siguiente pregunta de investigación: ¿Cómo influye el uso de semillas certificadas en el rendimiento del cultivo de arroz cáscara en las unidades agropecuarias del Perú durante el año 2024?

Para responder esta interrogante, la investigación utilizará la fuente de datos de la Encuesta Nacional Agropecuaria (ENA) 2024, diseñada por el Instituto Nacional de Estadística e Informática \citep{inei2024a}. Esta fuente es ideal por su representatividad nacional y diseño probabilístico, ya que permite un análisis integrado y sin sesgos al registrar variables críticas (producción, uso de semillas certificadas y factores de control).

Existe amplia evidencia de que las semillas certificadas aumentan el rendimiento agrícola. Sin embargo, dado que diversos factores exógenos —como el acceso a riego, el clima y las características territoriales— pueden confundir este impacto, resulta indispensable aislarlos mediante un análisis econométrico multivariado. En este contexto, el objetivo general de la presente investigación es determinar el efecto del uso de semillas certificadas sobre el rendimiento agrícola del cultivo de arroz cáscara en las unidades agropecuarias del Perú durante el año 2024.

Con esto, el estudio pretende generar evidencia empírica rigurosa que sirva de aporte para el diseño y optimización de políticas públicas agrícolas orientadas a la difusión tecnológica, el fortalecimiento de los sistemas de extensión agraria, la mejora de la productividad y la sostenibilidad de los productores de arroz en el país.

\section{Teoría, conceptos y antecedentes}

\subsection{Marco teórico}
Para analizar la relación entre el uso de semillas certificadas y el rendimiento del arroz cáscara, resulta necesario partir de la teoría de la producción, debido a que esta permite explicar cómo la calidad y combinación de los insumos empleados influyen sobre el producto obtenido por unidad de superficie. En la presente investigación, la semilla certificada es considerada un insumo tecnológico, mientras que el rendimiento agrícola constituye el resultado productivo que se busca explicar. \citet{just1979} sostienen que los insumos agrícolas no solo modifican el nivel esperado de producción, sino también la eficiencia y el riesgo asociados al proceso productivo. De manera complementaria, \citet{alston2021} señalan que la innovación tecnológica constituye uno de los principales determinantes del crecimiento de la productividad agrícola, debido a que permite incrementar la producción sin aumentar proporcionalmente el uso de tierra, trabajo o capital. En consecuencia, la calidad de los insumos puede generar diferencias relevantes en el rendimiento alcanzado por las unidades agropecuarias. En esta misma línea, \citet{li2024} concluyen que la tecnología agrícola contribuye a mejorar la eficiencia en el uso de los recursos y constituye un componente central de la intensificación sostenible. Asimismo, \citet{louwaars2021} destacan que la calidad genética, fisiológica y sanitaria de la semilla condiciona el establecimiento y desempeño posterior del cultivo. Bajo este enfoque, el empleo de semillas certificadas puede desplazar la función de producción hacia niveles superiores de rendimiento, al favorecer una mayor uniformidad y un mejor aprovechamiento del agua y de los demás insumos productivos. Por tanto, la teoría de la producción sustenta la hipótesis de que las unidades agropecuarias que utilizan semillas certificadas presentan un rendimiento del arroz cáscara superior al de aquellas que emplean semillas no certificadas, manteniendo constantes las demás condiciones incluidas en el modelo econométrico.

La teoría de la difusión de innovaciones resulta pertinente para la investigación porque el uso de semillas certificadas no depende únicamente de sus características técnicas, sino también de la decisión del productor de adoptarlas. En este sentido, esta teoría permite comprender por qué algunas unidades agropecuarias incorporan semillas certificadas dentro de su proceso productivo, mientras otras continúan utilizando semillas tradicionales o no certificadas. \citet{rogers2003} sostiene que la adopción de una innovación depende de la ventaja relativa percibida, su compatibilidad con las prácticas existentes, su complejidad y la posibilidad de observar o experimentar sus resultados. En el ámbito agrícola, \citet{macours2019} concluye que los productores no evalúan las innovaciones exclusivamente por su capacidad para elevar el rendimiento, sino también por sus efectos sobre el riesgo, el uso de agua, el trabajo requerido y la rentabilidad esperada. Por ello, una tecnología puede presentar ventajas agronómicas y, aun así, registrar una adopción limitada cuando sus beneficios no son claramente percibidos o existen restricciones económicas e informativas. Asimismo, \citet{okello2023} encontraron que los incentivos sociales y la circulación de información entre productores favorecen la difusión del conocimiento y la disposición a pagar por semillas certificadas. De manera complementaria, \citet{elbakali2023} concluyen que la extensión agraria, el acceso a información y las políticas públicas desempeñan un papel importante en la difusión de las innovaciones agrícolas. Estos resultados indican que la adopción de semillas certificadas responde tanto a sus beneficios productivos como al entorno institucional y social en el que toman decisiones los agricultores. En relación con el presente estudio, el uso de semillas certificadas representa una innovación tecnológica cuya adopción puede generar diferencias observables en el rendimiento del arroz cáscara. Por consiguiente, la teoría de la difusión de innovaciones complementa a la teoría de la producción, debido a que la primera explica la decisión de incorporar la tecnología, mientras que la segunda permite analizar su relación con los resultados productivos.

La inclusión de variables categóricas a nivel departamento es necesaria porque permite controlar la heterogeneidad espacial no observada y efectos fijos geográficos. Las regiones del Perú presentan marcadas diferencias en relación a la calidad de los suelos y la altitud \citep{escobal2002}. De acuerdo con la teoría de desarrollo agrícola, estas características influyen en los costos de transacción que enfrentan los productores y condicionan su acceso a los servicios de asistencia técnica como los canales de comercialización \citep{key2000}. Para el modelo econométrico la variable departamento es incorporada mediante variables dicotómicas (dummy) a fin de controlar las particularidades geográficas e institucionales.

Incorporar el sexo del productor como variable de control tiene el fin de considerar las desigualdades que pueden tener una influencia en el desempeño de la actividad agrícola. Estudios sostienen que las parcelas gestionadas por hombres muestran mayores rendimientos promedio frente a las gestionadas por mujeres, en principio por las limitaciones estructurales que enfrentan para acceder a recursos productivos esenciales como títulos de propiedad de la tierra, financiamiento formal, mercado de insumos y asistencia técnica \citep{aguilar2015}. En el presente estudio, esta variable de control se presenta mediante una variable dicotómica donde el valor de 1 representa a los productores hombres y el valor de 0 a las productoras mujeres.

El nivel de instrucción del productor constituye un elemento esencial del capital humano ya que influye en su capacidad para mejorar la eficiencia productiva y capacidad de absorción tecnológica. Los productores con mayores niveles educativos muestran una tasa de adopción de semillas certificadas significativamente más alta y una mayor eficiencia en el manejo de cultivos intensivos en conocimiento como el arroz \citep{ma2020}. La variable se define de una forma dicotómica, tomando el valor de uno si el jefe de la unidad agropecuaria posee educación superior (completa o incompleta) y cero si cuenta con educación básica o ninguna.

La edad representa la experiencia acumulada por el productor a lo largo de los años así como la etapa del ciclo laboral en la que se encuentra el agricultor. A más edad suele asociarse con mayor conocimiento pero también puede implicar menor capacidad física y una menor resistencia a asumir riesgos provocando una menor disposición a incorporar innovaciones tecnológicas y técnicas agrícolas más eficientes \citep{ouedraogo2021}.

\subsection{Antecedentes}
Diversas investigaciones han abordado la relación estructural entre el uso de semillas certificadas y el rendimiento agrícola, empleando diversas fuentes de microdatos y metodologías cuantitativas según el contexto. Estos estudios coinciden en que la calidad de la semilla es un determinante central del desempeño productivo, aunque su magnitud varía según condiciones agroecológicas y dotaciones de capital humano.

Un primer estudio realizado por \citet{takeshima2025} evaluó empíricamente la relación entre la producción local de semillas certificadas y su posterior adopción a nivel de parcela, cuantificando el efecto directo sobre el rendimiento físico de tres cultivos estratégicos (arroz, maíz y caupí), así como la heterogeneidad espacial por factores agroecológicos y socioeconómicos en Nigeria. Para tal fin, los autores emplearon microdatos longitudinales del panel de hogares agrícolas del \textit{General Household Survey-Panel} (GHS-Panel), emparejados geográficamente con registros de empresas procesadoras de semillas. En su estrategia empírica, aplicaron modelos de datos de panel con efectos fijos intra-parcela (\textit{panel fixed effects}) y variables instrumentales (VI/2SLS) para controlar la heterogeneidad inobservada y resolver la endogeneidad en la adopción varietal. Como resultados, demostraron un impacto positivo y altamente significativo sobre los rendimientos físicos de los cultivos, identificando un efecto marginal decreciente y heterogéneo que resulta marcadamente superior en parcelas ubicadas en zonas agroecológicas favorables con disponibilidad hídrica y productores con mayor nivel educativo. En relación con la presente investigación, este trabajo constituye el principal referente empírico al analizar específicamente el atributo de certificación oficial y justificar la inclusión de controles de estrés hídrico y climáticos para evitar sesgos por variable omitida.

De manera complementaria, un segundo estudio desarrollado por \citet{abdul2021} determinó cuantitativamente el impacto de la adopción de variedades mejoradas de arroz sobre el rendimiento agrícola y la eficiencia técnica en parcelas gestionadas por pequeños productores en Ghana. La base empírica consistió en microdatos primarios de corte transversal recolectados durante la campaña 2016 en 412 pequeños productores arroceros (286 adoptantes de variedades mejoradas y 126 no adoptantes). Para abordar la heterogeneidad tecnológica y la autoselección, desplegaron un enfoque econométrico multi-etapa que integró modelos Probit para los determinantes de adopción, emparejamiento por puntaje de propensión (\textit{Propensity Score Matching} - PSM) para mitigar el sesgo de selección observable, y Frontera de Producción Estocástica con análisis de meta-frontera para estimar brechas de eficiencia de manera insesgada. Las estimaciones revelaron un notable incremento y superioridad en la productividad física agrícola en comparación con los no adoptantes, al tiempo que la meta-frontera demostró una mayor cercanía a la frontera potencial de producción gracias a la resistencia genética ante tensiones bióticas y abióticas. Dicho hallazgo aporta evidencia sobre el salto productivo del germoplasma moderno y justifica metodológicamente la incorporación de controles sociodemográficos e hídricos en la regresión.

En el plano nacional, un tercer antecedente fundamental es el estudio de \citet{matallana2022}, quienes cuantificaron la rentabilidad financiera, los retornos netos de corto y largo plazo, así como las repercusiones ambientales derivadas de la adopción y sustitución por semilla certificada de arroz (\textit{Oryza sativa} L.) en los valles arroceros del Perú. Los investigadores recopilaron registros e información agropecuaria secundaria proveniente del Ministerio de Desarrollo Agrario y Riego (MIDAGRI), del INIA y de costos oficiales de producción en la costa norte (Lambayeque, La Libertad, Piura) y selva norte (San Martín, Amazonas). Su metodología consistió en un modelo de presupuesto parcial (\textit{partial budget analysis}) enfocado en la evaluación económico-financiera marginal, robustecido mediante simulaciones estocásticas de Montecarlo (10,000 iteraciones) para evaluar el riesgo en el Valor Actual Neto (VAN) y la Tasa Interna de Retorno (TIR) ante volatilidades climáticas y cambiarias. Los resultados proyectaron que el reemplazo de semilla común por certificada genera un incremento significativo en la rentabilidad neta del productor arrocero tanto en el corto como en el largo plazo, junto con una sustancial reducción del impacto ambiental motivada por el menor uso de plaguicidas. Este trabajo resulta relevante por compartir el ámbito geográfico y cultivo, aunque contrasta metodológicamente al usar presupuestos parciales con datos secundarios, frente a la estimación econométrica de rendimiento físico real mediante microdatos de la ENA 2024.

En conclusión, se adopta el estudio de \citet{takeshima2025} como la referencia metodológica y empírica central de la investigación, ya que este estudio aísla el impacto marginal del atributo de \textit{semilla certificada} sobre el rendimiento físico utilizando microdatos agropecuarios de alcance nacional (\textit{General Household Survey-Panel}), diseño muestral homólogo a la Encuesta Nacional Agropecuaria (ENA) del INEI. Por consiguiente, su marco conceptual fundamenta directamente la estimación econométrica por MCO con controles hídricos y territoriales en el contexto peruano.

\clearpage
\section{Metodología}

\subsection{Método y tipo de investigación}
La investigación adoptó una naturaleza cuantitativa, de tipo explicativo-correlacional, orientada a estimar el efecto del uso de semillas certificadas sobre el rendimiento agrícola del cultivo de arroz cáscara en las unidades agropecuarias del Perú. El diseño fue no experimental, de corte transversal y \textit{ex post facto}, debido a que las variables se observaron tal como se presentan en la realidad, sin manipulación, en un único corte temporal correspondiente al año 2024.

El estudio dispuso como técnica e instrumento de recolección de datos la revisión documental y la ficha de recopilación, respectivamente. La fuente de información fue la Encuesta Nacional Agropecuaria (ENA) 2024, ejecutada por el \citet{inei}. Esta encuesta cuenta con representatividad nacional y un diseño muestral probabilístico, además de constituir la fuente oficial para el estudio de la actividad agropecuaria en el país. Para efectos de la presente investigación, se emplearon dos capítulos de la ENA 2024: el Capítulo 200 (Partes A y B) [Módulo 1895], que registra la producción, superficie, prácticas de siembra y factores climáticos por parcela, y el Capítulo 1100 [Módulo 1911], que recoge las características sociodemográficas del productor/a agropecuario/a y su hogar. Ambos capítulos se integraron mediante un procedimiento de \textit{merge} (unión de base de datos) por medio de las llaves geográficas y de identificación de la unidad agropecuaria.

La población estuvo conformada por las unidades agropecuarias productoras de arroz cáscara en el Perú. Para la determinación de la muestra, se aplicaron los siguientes criterios de exclusión sobre la población: i) se conservaron únicamente las observaciones correspondientes al cultivo de arroz cáscara; ii) se restringió la muestra a unidades agropecuarias con una sola cosecha anual, con el propósito de reducir la heterogeneidad asociada a distintos calendarios de siembra; iii) se omitieron las observaciones a las que no les fue posible emparejarse entre ambos capítulos de la encuesta; iv) se excluyeron los registros con valores no positivos o perdidos en el rendimiento agrícola y en la variable de semillas certificadas (regresor principal), a razón de impedir estos la construcción del modelo; y v) se eliminó la observación correspondiente al departamento de Ayacucho por constituir un grupo singleton (una sola observación), dado que, de acuerdo con \citet{correia2015}, mantener grupos singleton en regresiones con efectos fijos anidados dentro de clusters puede sobrestimar la significancia estadística y conducir a inferencias incorrectas. Bajo la aplicación de los anteriores criterios, la muestra final del estudio quedó compuesta por 1,991 unidades agropecuarias.

\subsection{Especificación del modelo y definición de variables}
Dada la naturaleza continua de la variable dependiente, el rendimiento agrícola, se empleó un modelo de regresión lineal múltiple estimado por el método de Mínimos Cuadrados Ordinarios (MCO), aquella técnica econométrica que permite estimar los parámetros de regresión lineal minimizando la suma de los cuadrados de los residuos, asegurando bajo ciertos supuestos estimadores insesgados y eficientes \citep{wooldridge}. En su continuación, el rendimiento se transformó a su logaritmo natural (modelo log-lineal) con el objetivo de reducir la asimetría de su distribución y posibilitar que los coeficientes se interpreten como cambios porcentuales aproximados ante variaciones en las variables explicativas \citep{lozares1998}.

Por tanto, el modelo econométrico general se expresa de la siguiente manera:
\begin{equation}
	\ln(\text{Rend.}_i) = \beta_0 + \beta_1 \, \text{Semillas}_i + X_i \cdot \hat{\Gamma} + \varepsilon_i
\end{equation}

% --- TABLA 1: DEFINICIÓN DE VARIABLES (APARECERÁ AL INICIO DE PÁGINA 7) ---
\begin{table*}[t]
\raggedright
\footnotesize 
\caption{Definición y medición de las variables del modelo}
\label{tab:variables}
\begin{tabularx}{\textwidth}{@{} l l X c @{}}
\toprule
\textbf{Variable} & \textbf{Notación} & \textbf{Definición} & \textbf{Medición} \\
\midrule
Rendimiento agrícola ($\ln$) & \texttt{ln\_rendimiento} & Logaritmo natural del rendimiento del cultivo (producción en kg entre superficie cosechada en hectáreas) & \begin{tabular}[c]{@{}c@{}}Cuantitativa\\ Continua\end{tabular} \\
\addlinespace[0.08cm]
Semillas certificadas & \texttt{semillas\_certificadas} & \begin{tabular}[c]{@{}l@{}}0 = No certificada\\ 1 = Certificada\end{tabular} & Dicotómica \\
\addlinespace[0.08cm]
Fuente de agua & \texttt{fuente\_agua} & \begin{tabular}[c]{@{}l@{}}0 = Solo lluvia (secano)\\ 1 = Otra fuente (río, pozo, represa, etc.)\end{tabular} & Dicotómica \\
\addlinespace[0.08cm]
Sequía & \texttt{sequia} & \begin{tabular}[c]{@{}l@{}}0 = No reportó sequía\\ 1 = Reportó sequía como factor adverso\end{tabular} & Dicotómica \\
\addlinespace[0.08cm]
Lluvias a destiempo & \texttt{lluvias\_destiempo} & \begin{tabular}[c]{@{}l@{}}0 = No reportó lluvias fuera de temporada\\ 1 = Sí reportó\end{tabular} & Dicotómica \\
\addlinespace[0.08cm]
Plagas y enfermedades & \texttt{plagas\_enfermedades} & \begin{tabular}[c]{@{}l@{}}0 = No reportó plagas o enfermedades\\ 1 = Sí reportó\end{tabular} & Dicotómica \\
\addlinespace[0.08cm]
Otros factores climáticos & \texttt{otros\_factores} & \begin{tabular}[c]{@{}l@{}}0 = No reportó otros factores adversos\\ 1 = Sí reportó\end{tabular} & Dicotómica \\
\addlinespace[0.08cm]
Mujer productora & \texttt{mujer\_productora} & \begin{tabular}[c]{@{}l@{}}0 = Productor hombre\\ 1 = Productora mujer\end{tabular} & Dicotómica \\
\addlinespace[0.08cm]
Educación superior & \texttt{educacion\_superior} & \begin{tabular}[c]{@{}l@{}}0 = Educación básica o sin nivel\\ 1 = Educación superior ($\text{educ\_productor} \ge 7$)\end{tabular} & Dicotómica \\
\addlinespace[0.08cm]
Edad del productor & \texttt{edad\_productor} & Años de edad cumplidos del productor/a agropecuario & \begin{tabular}[c]{@{}c@{}}Cuantitativa\\ Continua\end{tabular} \\
\addlinespace[0.08cm]
Departamento & \texttt{departamento} & Departamento donde se ubica la unidad agropecuaria; efecto fijo regional (categoría base = San Martín) & \begin{tabular}[c]{@{}c@{}}Categórica\\ (16 cat.)\end{tabular} \\
\bottomrule
\end{tabularx}
\vspace{2pt}\par\noindent{\footnotesize \textit{Nota.} La variable dependiente es el logaritmo natural del rendimiento agrícola (kg/ha); la variable independiente principal ($X_1$) es el uso de semillas certificadas. Las demás son variables de control: disponibilidad hídrica, choques climáticos, características del productor/a (sexo, educación, edad) y efectos fijos departamentales.\par}
\end{table*}

Donde $X_i$ representa un vector de variables de control incorporadas con el objetivo de aislar el efecto del uso de semillas certificadas sobre el rendimiento agrícola y reducir el posible sesgo por omisión de variables relevantes. Dicho vector comprende variables asociadas a las condiciones productivas del cultivo, como el acceso a una fuente de agua distinta de la lluvia; factores climáticos y fitosanitarios que afectaron la producción durante la campaña agrícola (sequía, lluvias a destiempo, plagas y enfermedades, y otros factores); características sociodemográficas del productor (sexo, edad y nivel educativo); y la ubicación geográfica de la unidad agropecuaria, representada por el departamento donde se desarrolla la actividad agrícola.

A partir de la especificación presentada anteriormente, el modelo econométrico estimado incorpora explícitamente la variable explicativa principal y las variables de control seleccionadas, obteniéndose la siguiente especificación funcional desplegada:
\begin{equation}
\scriptsize
\begin{split}
    \ln(\text{Rendimiento})_i &= \beta_0 + \beta_1 \text{Semillas\_Certificadas}_i + \beta_2 \text{Fuente\_Agua}_i \\
    &\quad + \beta_3 \text{Sequía}_i + \beta_4 \text{Lluvias\_Destiempo}_i \\
    &\quad + \beta_5 \text{Plagas\_Enfermedades}_i + \beta_6 \text{Otros\_Factores}_i \\
    &\quad + \beta_7 \text{Mujer\_Productora}_i + \beta_8 \text{Educación\_Superior}_i \\
    &\quad + \beta_9 \text{Edad\_Productor}_i + \text{ib14.Departamento}_i + \varepsilon_i
\end{split}
\end{equation}

La \hyperref[tab:variables]{Tabla~\ref*{tab:variables}} resume formalmente la notación, definición conceptual y forma de medición de todas las variables empleadas en la estimación del modelo econométrico.

\subsection{Estadísticas descriptivas}
Tal como se reporta en la \hyperref[tab:descriptivas]{Tabla~\ref*{tab:descriptivas}}, el logaritmo natural del rendimiento agrícola ($\ln(\text{Rendimiento})$) presenta una media de $8.9589$ ($\text{DE} = 0.4752$), con valores mínimo y máximo de $7.0048$ y $9.6411$, respectivamente. Respecto a la variable independiente principal, el 78.36 \% de las parcelas de la muestra ($\text{Media} = 0.7836$) utilizó semillas certificadas durante la campaña agrícola 2024.

En cuanto a las variables de control, el 94.68 \% de las unidades agropecuarias registró una fuente de agua superficial o subterránea distinta a la lluvia exclusivamente ($\text{Media} = 0.9468$). Por el lado de los choques adversos durante la campaña, la mayor incidencia corresponde a plagas y enfermedades (22.19 \% de las parcelas), seguido por sequías (3.71 \%), otros factores adversos (2.01 \%) y lluvias a destiempo (1.51 \%). 

Finalmente, las características sociodemográficas revelan que la edad del productor agropecuario tiene un promedio de $55.597$ años ($\text{DE} = 13.236$, con un rango muestral de $17.000$ a $98.000$ años), mientras que la participación de mujeres productoras como conductoras principales es del 13.20 \% ($\text{Media} = 0.1320$) y la proporción de agricultores con educación superior se sitúa en 18.22 \% ($\text{Media} = 0.1822$).

% --- TABLA 2: ESTADÍSTICAS DESCRIPTIVAS ---
\begin{table}[ht!]
\raggedright
\scriptsize
\caption{Estadísticas descriptivas ($N=1,991$)}
\label{tab:descriptivas}
\begin{tabularx}{\columnwidth}{@{}Xcccc@{}}
\toprule
\textbf{Variable} & \textbf{Media} & \textbf{Desv. Est.} & \textbf{Mín.} & \textbf{Máx.} \\
\midrule
$\ln(\text{Rendimiento})$ & 8.9597 & 0.4742 & 7.0048 & 9.6411 \\
Semillas certificadas     & 0.7840 & 0.4116 & 0.0000 & 1.0000 \\
Fuente de agua            & 0.9473 & 0.2236 & 0.0000 & 1.0000 \\
Sequía                    & 0.0372 & 0.1892 & 0.0000 & 1.0000 \\
Lluvias a destiempo       & 0.0151 & 0.1219 & 0.0000 & 1.0000 \\
Plagas y enfermedades     & 0.2215 & 0.4154 & 0.0000 & 1.0000 \\
Otros factores            & 0.0201 & 0.1403 & 0.0000 & 1.0000 \\
Mujer productora          & 0.1321 & 0.3387 & 0.0000 & 1.0000 \\
Educación superior        & 0.1818 & 0.3858 & 0.0000 & 1.0000 \\
Edad del productor        & 55.6002 & 13.2389 & 17.0000 & 98.0000 \\
\bottomrule
\end{tabularx}
\vspace{2pt}\par\noindent{\scriptsize \textit{Nota.} Elaboración propia con base en microdatos de la ENA 2024 (INEI). Panel depurado de $N=1,991$ parcelas de arroz cáscara (sin el singleton de Ayacucho).\par}
\end{table}

\subsection{Resultados y estimación del modelo}
La \hyperref[tab:resultados]{Tabla~\ref*{tab:resultados}} reporta las estimaciones de los cinco modelos OLS anidados para el logaritmo natural del rendimiento agrícola ($\ln(\text{Rendimiento})$). Siguiendo el criterio metodológico de \citet{wooldridge}, la inclusión progresiva de bloques de control permite examinar la estabilidad paramétrica y el sesgo por omisión de variables sobre el coeficiente de interés ($\hat{\beta}_1$).

% --- TABLA 3: RESULTADOS DE LAS ESTIMACIONES OLS ---
\begin{table*}[t]
\raggedright 
\scriptsize 
\caption{Resultados progresivos de las estimaciones OLS para $\ln(\text{Rendimiento agrícola})$ }
\label{tab:resultados}
\begin{tabular*}{\textwidth}{@{\extracolsep{\fill}}lccccc@{}}
\toprule
\textbf{Variables Independientes} & \textbf{Modelo (1)} & \textbf{Modelo (2)} & \textbf{Modelo (3)} & \textbf{Modelo (4)} & \textbf{Modelo (5)} \\
\midrule
Semillas certificadas		& 0.2509*** & 0.0266 & 0.0202 & 0.0982*** & {0.0982***} \\
                            & (0.0253)  & (0.0226) & (0.0224) & (0.0234)  & {(0.0234)} \\
Fuente de agua              &           & 1.1527*** & 1.1311*** & 0.4858*** & {0.4858***} \\
                            &           & (0.0422) & (0.0422) & (0.0872)  & {(0.0871)} \\
Sequía                      &           & -0.0628  & -0.0593  & -0.0818*  & {-0.0818*} \\
                            &           & (0.0463) & (0.0459) & (0.0492)  & {(0.0492)} \\
Lluvias a destiempo         &           & -0.2819*** & -0.2874*** & -0.1968* & {-0.1968*} \\
                            &           & (0.0724) & (0.0718) & (0.1025)  & {(0.1025)} \\
Plagas y enfermedades       &           & -0.1253*** & -0.1307*** & -0.1323*** & {-0.1323***} \\
                            &           & (0.0210) & (0.0209) & (0.0200)  & {(0.0200)} \\
Otros factores climáticos   &           & -0.4151*** & -0.4079*** & -0.3629*** & {-0.3629***} \\
                            &           & (0.0617) & (0.0612) & (0.0625)  & {(0.0625)} \\
Mujer productora            &           &          & -0.0387  & -0.0733*** & {-0.0733***} \\
                            &           &          & (0.0254) & (0.0249)  & {(0.0249)} \\
Educación superior          &           &          & 0.1319*** & 0.0671*** & {0.0671***} \\
                            &           &          & (0.0226) & (0.0191)  & {(0.0191)} \\
Edad del productor          &           &          & 0.0014** & 0.0001  & {0.0001} \\
                            &           &          & (0.0007) & (0.0005)  & {(0.0005)} \\
Amazonas                    &           &          &          & 0.0738* (0.0443)   & 0.0738* (0.0443) \\
Áncash                      &           &          &          & 0.4436*** (0.0468) & 0.4436*** (0.0468) \\
Arequipa                    &           &          &          & 0.4987*** (0.0412) & 0.4987*** (0.0412) \\
Ayacucho                    &           &          &          & -0.7023*** (0.1037) & - \\
Cajamarca                   &           &          &          & 0.1531*** (0.0503) & 0.1531*** (0.0503) \\
Huánuco                     &           &          &          & -0.2049* (0.1085)  & -0.2049* (0.1085) \\
Junín                       &           &          &          & -0.8428*** (0.2387) & -0.8428*** (0.2387) \\
La Libertad                 &           &          &          & 0.3657*** (0.0411) & 0.3657*** (0.0411) \\
Lambayeque                  &           &          &          & 0.0975** (0.0406)  & 0.0975** (0.0406) \\
Loreto                      &           &          &          & -0.2290 (0.1453)   & -0.2290 (0.1453) \\
Madre de Dios               &           &          &          & -0.6204*** (0.1295) & -0.6204*** (0.1295) \\
Pasco                       &           &          &          & -0.5369*** (0.1465) & -0.5369*** (0.1465) \\
Piura                       &           &          &          & 0.0467 (0.0440)    & 0.0467 (0.0440) \\
San Martín (Base)           &           &          &          & 0.0000 (.)         & 0.0000 (.) \\
Tumbes                      &           &          &          & 0.1179* (0.0610)   & 0.1179* (0.0610) \\
Ucayali                     &           &          &          & -0.3421*** (0.0614) & -0.3421*** (0.0614) \\
\addlinespace[0.1cm]
Constante					& 8.7623*** & 7.8895*** & 7.8169*** & 8.2564*** & {8.2564***} \\
                            & (0.0224)  & (0.0386) & (0.0512) & (0.1046)  & {(0.1046)} \\
\midrule
Observaciones ($N$)         & 1,992     & 1,992    & 1,992    & 1,992    & {1,991} \\
$R^2$                       & 0.0473    & 0.3441   & 0.3564   & 0.5502   & {0.5480} \\
$R^2$ ajustado              & 0.0468    & 0.3421   & 0.3535   & 0.5447   & {0.5428} \\
\bottomrule
\end{tabular*}
\vspace{2pt}\par\noindent{\scriptsize \textit{Nota.} Errores estándar entre paréntesis (estimaciones robustas a heterocedasticidad de Huber-White en los Modelos 4 y 5). Nivel de significancia: *** $p<0{,}01$, ** $p<0{,}05$, * $p<0{,}1$.\par}
\end{table*}

\subsubsection{Estimación del modelo final}
El modelo empírico definitivo con efectos fijos departamentales (Modelo 5) alcanza un coeficiente de determinación de $R^2 = 0{,}5480$ ($R^2$ ajustado $0{,}5428$) y un estadístico global $F(23, 1967) = 76{,}32$ ($p < 0{,}0001$). 

Así, la ecuación empírica final ajustada por MCO se expresa de forma sintética con números y vector de controles ($X_i \cdot \hat{\Gamma}$) de la siguiente manera:

\begin{equation}
\ln(\text{Rend.}_i) = 8{,}2564 + 0{,}0982 \, \text{Semillas}_i + X_i \cdot \hat{\Gamma} + \varepsilon_i
\end{equation}
donde $\text{Semillas}_i$ es la variable dicotómica de adopción ($X_{1i}$) y $X_i \cdot \hat{\Gamma}$ representa el vector estimado que engloba tanto las variables observables de control a nivel individual como los efectos fijos espaciales departamentales ($\hat{\mu}_d$).

Asimismo, al desplegar explícitamente en texto cada variable con su coeficiente numérico estimado en el Modelo 5, la ecuación ajustada completa se expresa como:

\begin{equation}
\begin{split}
\ln(\text{Rend.}_i) &= 8{,}2564 + 0{,}0982 \, \text{Semillas}_i \\
&\quad + 0{,}4858 \, \text{Agua}_i - 0{,}0818 \, \text{Sequía}_i \\
&\quad - 0{,}1968 \, \text{Lluvias}_i - 0{,}1323 \, \text{Plagas}_i \\
&\quad - 0{,}3629 \, \text{Otros\_Clima}_i - 0{,}0733 \, \text{Mujer}_i \\
&\quad + 0{,}0671 \, \text{Educ.}_i + 0{,}0001 \, \text{Edad}_i \\
&\quad + \text{ib14.Departamento}_i
\end{split}
\end{equation}

En correspondencia con el punto 3.2, el contraste entre los Modelos 4 ($N=1{,}992$) y 5 ($N=1{,}991$) de la \hyperref[tab:resultados]{Tabla~\ref*{tab:resultados}} ilustra nítidamente el efecto de excluir la observación solitaria (\textit{singleton}) de Ayacucho. Si bien las pendientes se mantienen invariantes ($\hat{\beta}_1 = 0{,}0982$), depurar este grupo genera una doble mejora estructural: i) el $R^2$ desciende de $0{,}5502$ a $0{,}5480$, eliminando el sobreajuste artificial del 100\% de su intercepto propio y revelando un poder explicativo más real; y ii) los grados de libertad en \texttt{vce(robust)} se corrigen con exactitud, afinando la precisión inferencial, como se aprecia al contrastar el error estándar robusto de la fuente de agua: $0{,}0872$ en el Modelo 4 frente a $0{,}0871$ en el Modelo 5 al reportarse a cuatro decimales ($0{,}08717$ vs. $0{,}08714$ a cinco decimales en los registros unificados).

\subsubsection{Impacto de la semilla certificada}
En el modelo final (Modelo 5), el resultado empírico evidencia que la adopción de semilla certificada ejerce un impacto positivo y altamente significativo sobre la productividad agrícola ($\hat{\beta}_1 = 0{,}0982$, $p < 0{,}001$ y $t = 4{,}19$). Al multiplicar el coeficiente por cien para su interpretación directa, se demuestra que el uso de semilla certificada incrementa el rendimiento del arroz cáscara en $+9{,}82$\% en comparación con la semilla común no certificada, manteniendo constantes las condiciones de riego, los factores climáticos, el perfil demográfico del productor y los efectos territoriales.

\subsubsection{Controles y efectos fijos}
En el vector de controles ($X_i \cdot \hat{\Gamma}$), la aproximación directa ($\hat{\beta} \times 100$) en el orden de la regresión indica que: i) la fuente de agua aumenta el rendimiento en $+48{,}58$\% ($p < 0{,}001$); ii) los choques climáticos reducen la cosecha: sequía en $-8{,}18$\% ($p < 0{,}10$), lluvias a destiempo en $-19{,}68$\% ($p < 0{,}10$), plagas y enfermedades en $-13{,}23$\% ($p < 0{,}001$) y otros factores climáticos en $-36{,}29$\% ($p < 0{,}001$); y iii) en el perfil demográfico, la mujer productora presenta una brecha de $-7{,}33$\% ($p < 0{,}001$), la educación superior aporta $+6{,}71$\% ($p < 0{,}001$) y la edad resulta neutra ($+0{,}01$\%, $p = 0{,}807$).

Asimismo, los efectos fijos territoriales miden las diferencias en comparación directa con el departamento base (San Martín, Base $= 0{,}00$). En orden de estimación, respecto a San Martín: Amazonas obtiene $+7{,}38$\% ($p < 0{,}10$), Áncash $+44{,}36$\% ($p < 0{,}001$), Arequipa lidera con $+49{,}87$\% ($p < 0{,}001$) y Cajamarca $+15{,}31$\% ($p < 0{,}001$). Frente a la base, Huánuco y Junín muestran penalizaciones de $-20{,}49$\% ($p < 0{,}10$) y $-84{,}28$\% ($p < 0{,}001$), mientras que La Libertad y Lambayeque superan a San Martín con $+36{,}57$\% ($p < 0{,}001$) y $+9{,}75$\% ($p < 0{,}05$). Finalmente, respecto a San Martín: Loreto registra $-22{,}90$\% ($p > 0{,}10$), Madre de Dios y Pasco caen $-62{,}04$\% y $-53{,}69$\% ($p < 0{,}001$), Piura es similar a la base ($+4{,}67$\%, $p > 0{,}10$), Tumbes obtiene $+11{,}79$\% ($p < 0{,}10$) y Ucayali cae $-34{,}21$\% ($p < 0{,}001$).

\subsubsection{Validación del modelo}
Las pruebas post-estimación corroboran la solidez estructural del modelo: i) el VIF general es $2{,}06$ ($1{,}39$ en semilla $X_{1i}$), descartando multicolinealidad; ii) la prueba de Breusch-Pagan ($\chi^2 = 108{,}20$, $p < 0{,}0001$) rechaza homocedasticidad, validando errores robustos de Huber-White (\texttt{vce(robust)}) para inferencia insesgada en estadísticos $t$ y $F$; iii) la estabilidad de la semilla ($0{,}2509$ en Modelo 1 a $0{,}0982$ en Modelo 5) descarta sesgos por variable omitida; y iv) el test de significancia conjunta global ratifica su alto poder explicativo ($F(23, 1967) = 76{,}32$, $p < 0{,}0001$). Asimismo, se asume que la adopción de semilla ($X_{1i}$) conlleva cierta endogeneidad por autoselección (habilidad inobservada); por ende, el incremento estimado identifica una asociación robusta de alta relevancia agropecuaria.

\subsubsection{Discusión}
A partir de la rigurosidad estadística del modelo econométrico depurado y de las pruebas diagnósticas, se confirma empíricamente la hipótesis general de la investigación: el uso de semilla certificada ejerce un efecto positivo y estadísticamente significativo sobre el rendimiento físico del arroz cáscara en el Perú ($\hat{\beta}_1 = 0{,}0982$, $p < 0{,}001$). Este diferencial del $+9{,}82$\% ratifica la superioridad productiva del germoplasma formal frente a la semilla común reciclada, consolidándolo como un insumo estratégico para la agricultura nacional \citep{fao2019}.

Al contrastar este hallazgo con la literatura referencial, la evidencia obtenida en la ENA 2024 resulta plenamente consistente con lo reportado por \citet{takeshima2025} en Nigeria (\textit{General Household Survey-Panel}). En ambos estudios se demuestra que aislar estrictamente el atributo de \textit{certificación oficial} arroja un retorno marginal positivo sobre el rendimiento en parcela, y que dicha productividad depende estructuralmente del acceso al riego y del control por heterogeneidad agroecológica. La notable estabilidad de la pendiente en el modelo peruano consolida empíricamente la validez de los hallazgos de Takeshima et al. (2025) en encuestas agropecuarias probabilísticas homólogas.

\section{Conclusiones}
Utilizando microdatos probabilísticos de la Encuesta Nacional Agropecuaria (ENA) 2024, evaluamos el efecto del uso de semillas certificadas sobre el rendimiento agrícola a nivel de parcela en los productores de arroz cáscara del Perú, así como su heterogeneidad en función de factores agroecológicos, climáticos y sociodemográficos. Aplicamos modelos de regresión múltiple por Mínimos Cuadrados Ordinarios (MCO) con errores estándar robustos a heterocedasticidad de Huber-White y efectos fijos departamentales para controlar la heterogeneidad espacial no observada. Encontramos evidencia robusta de que la adopción de semilla certificada genera un impacto productivo positivo y altamente significativo ($p < 0{,}001$), incrementando el rendimiento por hectárea en $+9{,}82$\% frente al uso de semilla común o reciclada.

Asimismo, los resultados evidencian que el rendimiento arrocero es altamente heterogéneo y depende estructuralmente de las condiciones de la unidad agropecuaria. El acceso a riego tecnificado se consolida como el determinante individual con mayor impulso en la productividad ($+48{,}58$\%), mientras que los choques exógenos adversos —como plagas y anomalías hídricas— imponen penalidades severas que oscilan entre $-8{,}18$\% y $-36{,}29$\%. La incorporación de estos controles agroecológicos eleva el poder explicativo del modelo de un $4{,}73$\% a un sólido $54{,}80$\% ($R^2 = 0{,}5480$, $F = 76{,}32$), mostrando que la magnitud del coeficiente de semilla certificada se mantiene excepcionalmente estable ante la inclusión de covariables, lo que descarta sesgos graves por variable omitida.

En términos de política agraria, nuestros hallazgos subrayan la importancia de focalizar esfuerzos territoriales para fortalecer los sistemas de extensión y la absorción varietal. Se recomienda priorizar esquemas de subsidio temporal o crédito agrícola en parcelas de pequeña escala donde los costos de certificación representan una barrera de entrada, articulando estas intervenciones con la expansión de infraestructura de riego para maximizar el retorno físico en el campo. Finalmente, si bien la estimación lineal identifica una relación productiva consistente, investigaciones futuras podrán abordar la potencial endogeneidad por autoselección varietal mediante variables instrumentales o diseños experimentales.

% --- SECCIÓN DE REFERENCIAS APA 7 (NATBIB VINCULADO) ---
\vspace{-4pt}
\begin{thebibliography}{45}
\footnotesize
\setlength{\itemsep}{0pt}
\setlength{\parskip}{0pt}

\bibitem[Abdul-Rahaman et al.(2021)]{abdul2021} 
Abdul-Rahaman, A., Issahaku, G., \& Zereyesus, Y. A. (2021). Improved rice variety adoption and farm production efficiency: Accounting for unobservable selection bias and technology gaps among smallholder farmers in Ghana. \textit{Technology in Society}, 64, 101471. \url{https://doi.org/10.1016/j.techsoc.2020.101471}

\bibitem[Aguilar et al.(2015)]{aguilar2015} 
Aguilar, A., Carranza, E., Goldstein, M., Kilic, T., \& Oseni, G. (2015). Descomposición de las diferencias de género en la productividad agrícola en Etiopía. \textit{Agricultural Economics}, 46, 311--334. \url{https://doi.org/10.1111/agec.12167}

\bibitem[Alston et al.(2021)]{alston2021} 
Alston, J. M., Pardey, P. G., \& Rao, X. (2021). The economics of agricultural innovation. En C. B. Barrett \& D. R. Just (Eds.), \textit{Handbook of Agricultural Economics} (Vol. 5, pp. 3895--3980). Elsevier. \url{https://doi.org/10.1016/bs.hesagr.2021.10.001}

\bibitem[Battese et al.(2004)]{battese2004} 
Battese, G. E., Rao, D. S. P., \& O'Donnell, C. J. (2004). A metafrontier production function for estimation of technical efficiencies and technology gaps for firms operating under different technologies. \textit{Journal of Productivity Analysis}, 21, 91--103.

\bibitem[Correia(2015)]{correia2015} 
Correia, S. (2015). Singletons, cluster-robust standard errors and fixed effects: A bad mix. \textit{Technical Note, Duke University}, 7(9), 1--7. \url{https://scorreia.com/research/singletons.pdf}

\bibitem[Dewi et al.(2026)]{dewi2026} 
Dewi, Y. A., Bahru, B. A., \& Zeller, M. (2026). Revealing the effect of agricultural interventions on technical efficiency: Empirical evidence from Indonesian smallholder rice farmers. \textit{Agricultural and Food Economics}, 14, 16. \url{https://doi.org/10.1186/s40100-026-00454-1}

\bibitem[El Bakali et al.(2023)]{elbakali2023} 
El Bakali, I., Brouziyne, Y., Ait El Mekki, A., Maatala, N., \& Harbouze, R. (2024). The impact of policies on the diffusion of agricultural innovations: Systematic review on evaluation approaches. \textit{Outlook on Agriculture}, 53(1), 3--22. \url{https://doi.org/10.1177/00307270231215837}

\bibitem[Escobal y Ponce(2002)]{escobal2002} 
Escobal, J., \& Ponce, C. (2002). \textit{El beneficio de los caminos rurales en el Perú}. Grupo de Análisis para el Desarrollo (GRADE). \url{https://biblioteca.clacso.edu.ar/Peru/grade/20100511023603/ddt40ES.pdf}

\bibitem[Evenson y Gollin(2003)]{evenson2003} 
Evenson, R. E., \& Gollin, D. (2003). Assessing the impact of the Green Revolution, 1960 to 2000. \textit{Science}, 300(5620), 758--762. \url{https://doi.org/10.1126/science.1078710}

\bibitem[FAO(2011)]{fao2011} 
FAO. (2011). \textit{The state of the world's land and water resources for food and agriculture (SOLAW): Managing systems at risk}. Food and Agriculture Organization of the United Nations; Earthscan. \url{https://www.fao.org/3/i1688e/i1688e.pdf}

\bibitem[FAO(2016)]{fao2016} 
FAO. (2016). \textit{El Estado mundial de la agricultura y la alimentación 2016}. Food and Agriculture Organization of the United Nations. \url{https://www.fao.org/agrifood-economics/publications/detail/es/c/1179968/}

\bibitem[FAO y AfricaSeeds(2019)]{fao2019} 
FAO y AfricaSeeds. (2019). \textit{Materiales para capacitación en semillas. Módulo 3: Control de calidad y certificación de semillas}. Organización de las Naciones Unidas para la Alimentación y la Agricultura. \url{https://www.fao.org/3/ca1492es/CA1492ES.pdf}

\bibitem[FAO(2023)]{fao2023} 
FAO. (2023). \textit{FAOSTAT: Crops and livestock products. Metadata}. Food and Agriculture Organization of the United Nations. \url{https://www.fao.org/faostat/}

\bibitem[Feder et al.(1985)]{feder1985} 
Feder, G., Just, R. E., \& Zilberman, D. (1985). Adoption of agricultural innovations in developing countries: A survey. \textit{Economic Development and Cultural Change}, 33(2), 255--298. \url{https://doi.org/10.1086/451461}

\bibitem[Huang et al.(2014)]{huang2014} 
Huang, C. J., Huang, T. H., \& Liu, N. H. (2014). A new approach to estimating the metafrontier production function based on a stochastic frontier framework. \textit{Journal of Productivity Analysis}, 42, 241--254.

\bibitem[INEI(2024)]{inei} 
Instituto Nacional de Estadística e Informática. (2024). \textit{Encuesta Nacional Agropecuaria (ENA) 2024}. INEI.

\bibitem[INEI(2024a)]{inei2024a} 
Instituto Nacional de Estadística e Informática. (2024a). \textit{Encuesta Nacional Agropecuaria (ENA) 2024: Ficha técnica y metodología}. INEI. \url{https://www.gob.pe/inei}

\bibitem[INEI(2024b)]{inei2024b} 
Instituto Nacional de Estadística e Informática. (2024b). \textit{Encuesta Nacional Agropecuaria (ENA) 2024: Manual del encuestador e instrumentos de recolección}. INEI. \url{https://www.datosabiertos.gob.pe/dataset/encuesta-nacional-agropecuaria-ena-2024-instituto-nacional-de-estadistica-e-informatica-inei}

\bibitem[Just y Pope(1979)]{just1979} 
Just, R. E., \& Pope, R. D. (1979). Production function estimation and related risk considerations. \textit{American Journal of Agricultural Economics}, 61(2), 276--284. \url{https://doi.org/10.2307/1239732}

\bibitem[Key et al.(2000)]{key2000} 
Key, N., Sadoulet, E., \& Janvry, A. D. (2000). Costos de transacción y respuesta de la oferta de los hogares agrícolas. \textit{American Journal of Agricultural Economics}, 82, 245--259. \url{https://doi.org/10.1111/0002-9092.00022}

\bibitem[Li et al.(2024)]{li2024} 
Li, Y., Herzog, F., Levers, C., et al. (2024). Agricultural technology as a driver of sustainable intensification: Insights from the diffusion and focus of patents. \textit{Agronomy for Sustainable Development}, 44, Article 14. \url{https://doi.org/10.1007/s13593-024-00949-5}

\bibitem[Louwaars y De Jonge(2021)]{louwaars2021} 
Louwaars, N. P., \& De Jonge, B. (2021). Regulating seeds---A challenging task. \textit{Agronomy}, 11(11), 2324. \url{https://doi.org/10.3390/agronomy11112324}

\bibitem[Lozares et al.(1998)]{lozares1998} 
Lozares, C., López Roldán, P., \& Borrás, V. (1998). La complementariedad del log-lineal y del análisis de correspondencias en la elaboración y el análisis de tipología. \textit{Papers: Revista de Sociología}, 55, 79--93. \url{https://doi.org/10.5565/rev/papers.1933}

\bibitem[Ma y Wang(2020)]{ma2020} 
Ma, W., \& Wang, X. (2020). Internet use, sustainable agricultural practices and rural incomes: Evidence from China. \textit{Australian Journal of Agricultural and Resource Economics}, 64(4), 1081--1112. \url{https://ideas.repec.org/a/ags/aareaj/342932.html}

\bibitem[Macours(2019)]{macours2019} 
Macours, K. (2019). Farmers' demand and the traits and diffusion of agricultural innovations in developing countries. \textit{Annual Review of Resource Economics}, 11, 483--499. \url{https://doi.org/10.1146/annurev-resource-100518-094045}

\bibitem[Matallana et al.(2022)]{matallana2022} 
Matallana, R., Minaya, C., Vásquez Quispe, C., Barrientos, N., Duárez, A., Cusi, M., \& Velarde, S. (2022). Beneficios económicos y ambientales de la semilla certificada de arroz (Oryza sativa) en la costa norte y selva norte del Perú. \textit{Natura@economía}, 7(1), 63--80. \url{https://doi.org/10.21704/ne.v7i1.2111}

\bibitem[O'Donnell et al.(2008)]{odonnell2008} 
O'Donnell, C. J., Rao, D. S. P., \& Battese, G. E. (2008). Metafrontier frameworks for the study of firm-level efficiencies and technology ratios. \textit{Empirical Economics}, 34, 231--255.

\bibitem[Okello et al.(2023)]{okello2023} 
Okello, J. J., Shikuku, K. M., Lagerkvist, C. J., Rommel, J., Jogo, W., Ojwang, S., Namanda, S., \& Elungat, J. (2023). Social incentives as nudges for agricultural knowledge diffusion and willingness to pay for certified seeds: Experimental evidence from Uganda. \textit{Food Policy}, 120, 102506. \url{https://doi.org/10.1016/j.foodpol.2023.102506}

\bibitem[Ouédraogo et al.(2021)]{ouedraogo2021} 
Ouédraogo, M., Houessionon, P., Zougmoré, R. B., \& Partey, S. T. (2021). Factors influencing the adoption of climate-smart agricultural practices in rural Burkina Faso. \textit{Environmental Management}, 67(4), 704--718. \url{https://doi.org/10.18697/ajfand.125.23400}

\bibitem[Rogers(2003)]{rogers2003} 
Rogers, E. M. (2003). \textit{Diffusion of innovations} (5th ed.). Free Press.

\bibitem[Schlenker y Roberts(2009)]{schlenker2009} 
Schlenker, W., \& Roberts, M. J. (2009). Nonlinear temperature effects indicate severe damages to US crop yields under climate change. \textit{Proceedings of the National Academy of Sciences}, 106(37), 15594--15598. \url{https://doi.org/10.1073/pnas.0906865106}

\bibitem[Takeshima et al.(2025)]{takeshima2025} 
Takeshima, H., Ragasa, C., Bamiwoye, T., Andam, K. S., Spielman, D. J., Edeh, H. O., et al. (2025). Seed certification, certified seeds use and yield outcomes in Nigeria: Insights from nationally-representative farm panel data and seed company location data. \textit{Agricultural Systems}, 224, 104268. \url{https://doi.org/10.1016/j.agsy.2025.104268}

\bibitem[Wooldridge(2010)]{wooldridge} 
Wooldridge, J. M. (2010). \textit{Introducción a la econometría: Un enfoque moderno} (4.\textsuperscript{a} ed.). Cengage Learning.

\end{thebibliography}

\end{document}